% !PRINTED 361
\weberpage{361}
%<<CONT
then there follows from this the identity
\[
x_1 y_1 (x_2 y_2 - x_3 y_3) = (u_1 - u_2)(u_1 + u_2),
\]
and from this, as above,
\[
\begin{aligned}
u_1 - u_2 &= h x_1 y_1\\
u_1 + u_2 &= \frac{x_2 y_2 - x_3 y_3}{h},
\end{aligned}
\]
in which $h$ denotes a constant; hence
\[
2u_1 = h x_1 y_1 + \frac{x_2 y_2 - x_3 y_3}{h},
\]
and accordingly
\weq{7}{4f = 4 x_1 y_1 x_2 y_2 - \left(h x_1 y_1 + \frac{x_2 y_2 - x_3 y_3}{h}\right)^2.}
We can give this form of the equation a more elegant shape if we replace
\[
h x_1,\quad \frac{x_2}{h},\quad \frac{x_3}{h}
\]
by
\[
x_1,\quad x_2,\quad x_3
\]
respectively. Then the new $x_1, x_2, x_3$, set equal to zero, denote the same lines as the original ones, since the two indeed differ only by a constant factor, and from (7) there results:
\weq{8}{\begin{aligned}
-4f &= (x_1 y_1 + x_2 y_2 - x_3 y_3)^2 - 4 x_1 y_1 x_2 y_2\\
&= x_1^2 y_1^2 + x_2^2 y_2^2 + x_3^2 y_3^2\\
&\quad - 2 x_2 y_2 x_3 y_3 - 2 x_1 y_1 x_2 y_2 - 2 x_1 y_1 x_3 y_3,
\end{aligned}}
and the equation $f = 0$ can also be represented in the elegant irrational form
\[
x_1 y_1 + x_2 y_2 - x_3 y_3 = -2\sqrt{x_1 y_1 x_2 y_2}
\]
or
\weq{9}{\sqrt{x_1 y_1} + \sqrt{x_2 y_2} + \sqrt{x_3 y_3} = 0}
From (9) we can again derive the form (8) backwards. But since (9) is entirely symmetric, we can interchange the three pairs, and obtain, e.g., also
\[
4 x_2 y_2 x_3 y_3 = (-x_1 y_1 + x_2 y_2 + x_3 y_3)^2,
\]
from which it is to be seen that the points of contact of $x_2 y_2 x_3 y_3$ also lie on a conic section, and that in the complex which one obtains from $x_2 y_2$, not only the pair $x_1 y_1$, but
%CONT>>
